\documentclass[11pt]{article}
% Shared print layout. The web's provisional design and MathJax are independent.
\usepackage[utf8]{inputenc}
\usepackage[T1]{fontenc}
\usepackage[spanish,es-nodecimaldot]{babel}
\usepackage{lmodern}
\usepackage{amsmath,amssymb}
\usepackage{graphicx}
\usepackage{geometry}
\geometry{a4paper,left=20mm,right=20mm,top=25mm,bottom=22mm,headheight=16pt,headsep=8mm}
\usepackage{fancyhdr,lastpage}
\usepackage{enumitem}
\usepackage{xcolor}
\usepackage{pgfplots}
\pgfplotsset{compat=1.18}
\usepgfplotslibrary{fillbetween}
\usepackage[hidelinks,pdfencoding=auto]{hyperref}
\hypersetup{pdflang=es-ES,pdfauthor={},pdfsubject={Matemáticas Aplicadas a las Ciencias Sociales II}}
\pagestyle{fancy}
\fancyhf{}
\fancyhead[L]{\small Matemáticas Aplicadas a las Ciencias Sociales II}
\fancyhead[R]{\small 2.\textsuperscript{o} Bachillerato}
\fancyfoot[L]{\footnotesize Probabilidad II | Edición revisada 1.0}
\fancyfoot[R]{\footnotesize Página \thepage\ de \pageref{LastPage}}
\renewcommand{\headrulewidth}{0.4pt}
\setlength{\parindent}{0pt}
\setlength{\parskip}{0.5em}
\setlist{itemsep=0.3em,topsep=0.3em}
\newcommand{\apartado}[1]{\subsection*{#1}}
\pgfplotsset{
  revision/.style={
    width=0.94\linewidth,height=7.5cm,
    axis lines=middle,grid=major,grid style={gray!18},
    tick label style={font=\small},label style={font=\small},
    xlabel={$x$},ylabel={$y$},samples=160,
    legend style={at={(0.5,-0.16)},anchor=north,draw=none,font=\small},
    clip=true
  }
}

\hypersetup{pdftitle={Probabilidad II - Soluciones revisadas}}
\begin{document}
\renewcommand{\arraystretch}{2.2}
\section*{Ejercicio 1: sucesos y condicionamiento}
Soluciones revisadas a partir de F017. Se conservan todos los datos y los cuatro ejercicios, aunque solo dos de las tres opcionales se evalúan.

La tabla de probabilidades conjuntas completa permite comprobar los marginales y los complementos:
\[
\begin{array}{c|cc|c}
 & B & B^c & \text{Total}\\\hline
 A & \dfrac{3}{20} & \dfrac{1}{4} & \dfrac{2}{5}\\
 A^c & \dfrac{7}{20} & \dfrac{1}{4} & \dfrac{3}{5}\\\hline
 \text{Total} & \dfrac{1}{2} & \dfrac{1}{2} & 1
\end{array}
\]
\begin{enumerate}[label=(\alph*)]
\item $P(A\cap B)=P(A)-P(A\cap B^c)=\boxed{\frac{3}{20}}$.
\item $P(A^c\cap B)=P(B)-P(A\cap B)=\boxed{\frac{7}{20}}$.
\item $P(A\mid B)=P(A\cap B)/P(B)=\boxed{\frac{3}{10}}$, con $P(B)>0$.
\item $P(A)P(B)=\frac{1}{5}\ne \frac{3}{20}=P(A\cap B)$: no son independientes.
\end{enumerate}
\clearpage\section*{Ejercicio 2: lectura de diarios}
\begin{enumerate}[label=(\alph*)]
\item Por inclusión-exclusión, $P(A\cup B)=P(A)+P(B)-P(A\cap B)=\frac{9}{20}$. Por complemento, $P((A\cup B)^c)=\boxed{\frac{11}{20}}$, es decir, 55\%.
\item Se condiciona entre quienes no leen B; no entre toda la población. $P(B^c)=\frac{41}{50}>0$ y $P(A\cap B^c)=\frac{27}{100}$. Por tanto,\[P(A\mid B^c)=\frac{P(A\cap B^c)}{P(B^c)}=\boxed{\frac{27}{82}}\approx 0,3293.\]
\end{enumerate}\clearpage\section*{Ejercicio 3: alumnado y aprobados}
La selección es uniforme entre las 800 personas matriculadas. Sea $S$ el suceso aprobar.
\[\begin{array}{c|rr|r}&S&S^c&\text{Total}\\\hline\text{Hombres}&208&112&320\\\text{Mujeres}&360&120&480\\\hline\text{Total}&568&232&800\end{array}\]
\begin{enumerate}[label=(\alph*)]\item Por recuento directo,\[P(S)=\frac{568}{800}=\boxed{\frac{71}{100}}=0,7100.\]
La probabilidad total da el mismo resultado: $P(S)=P(H)P(S\mid H)+P(M)P(S\mid M)$.
\item Entre quienes aprueban, $P(M\mid S)=\boxed{\frac{45}{71}}\approx 0,6338$. Bayes confirma el cociente: $P(M)P(S\mid M)/P(S)$.
\end{enumerate}\clearpage\section*{Ejercicio 4: secaderos y Reserva}
Se interpreta la elección como uniforme entre todos los jamones, conforme al cálculo original. No se elige primero un secadero con probabilidad $1/2$. Si esa fuera la intención, habría que precisar el enunciado y cambiar el modelo. Sea $R$ el suceso ser de Reserva.
\[\begin{array}{c|cc|c}&R&R^c&\text{Total}\\\hline A&\dfrac{9}{50}&\dfrac{21}{50}&\dfrac{3}{5}\\ B&\dfrac{9}{25}&\dfrac{1}{25}&\dfrac{2}{5}\\\hline\text{Total}&\dfrac{27}{50}&\dfrac{23}{50}&1\end{array}\]
\begin{enumerate}[label=(\alph*)]\item Por probabilidad total,\[P(R^c)=(1-P(B))(1-P(R\mid A))+P(B)(1-P(R\mid B))=\boxed{\frac{23}{50}}\approx 0,4600.\]
\item Como $P(R)=\frac{27}{50}>0$, Bayes da\[P(A\mid R)=\frac{P(A)P(R\mid A)}{P(R)}=\boxed{\frac{1}{3}}\approx 0,3333.\]
\end{enumerate}
Se conservan las fracciones exactas y se redondea solo el resultado final.
\end{document}
